\section{Cryptography}

Cryptography is used to provide security properties such as
\textbf{confidentiality}, \textbf{integrity}, and \textbf{authentication}.


\paragraph{What cryptography is not}
Cryptography should not be confused with \emph{steganography}, which
hides the existence of a message, or with simple \emph{encoding}, which
only changes the representation of information.

\subsection{Symmetric-Key Cryptography}

In symmetric-key cryptography, the sender and the receiver share the
same secret key. The same key is used for encryption and decryption.

A symmetric cipher consists of two efficient algorithms:

\[
E : K \times M \rightarrow C
\]

\[
D : K \times C \rightarrow M
\]

where $K$ is the set of keys, $M$ the set of messages, and $C$ the set
of ciphertexts. Correctness means that decrypting an encrypted message
with the same key gives back the original message:

\[
\forall m \in M, \forall k \in K :
D(k, E(k,m)) = m
\]

The main difficulty is that the key must be shared secretly between the
sender and the receiver.

\paragraph{Perfect Security}

A cipher is perfectly secure if the ciphertext does not reveal any
information about the plaintext, except possibly its length. More
formally, encryptions of any two messages should have the same
distribution over all possible keys.

Perfect security is very strong, but it requires the key length to be at
least as long as the message length. This is usually impractical.

\paragraph{One-Time Pad}

The One-Time Pad (OTP) is a perfectly secure cipher. It uses a random key
$k$ of the same length as the message $m$:

\[
E(k,m) = k \oplus m
\]

\[
D(k,c) = k \oplus c
\]

where $\oplus$ is the bitwise XOR operation.

Decryption works because:

\[
D(k,E(k,m)) = k \oplus (k \oplus m) = m
\]

The main problem is that the key must be random, as long as the message,
and used only once.

\paragraph{Attack on Reused One-Time Pad Keys}

The OTP becomes insecure if the same key is reused for two messages:

\[
c_1 = m_1 \oplus k
\]

\[
c_2 = m_2 \oplus k
\]

An attacker can compute:

\[
c_1 \oplus c_2 = m_1 \oplus m_2
\]

The key disappears. Since human languages are redundant, an attacker can
often guess parts of one message and recover parts of the other. This is
why reusing a one-time pad key completely breaks the security.

\paragraph{Stream Ciphers}

A stream cipher tries to make the One-Time Pad practical. Instead of
sharing a key as long as the message, both parties share a short secret
seed $s$. A pseudo-random generator then expands this seed into a long
keystream:

\[
k_1, k_2, \ldots = PRG(s)
\]

The message is then encrypted by XORing it with the generated keystream.
The security depends on the quality of the pseudo-random generator and
on the secrecy of the seed.

\paragraph{Initialization Vector}

To avoid generating the same keystream twice, stream ciphers use an
\emph{Initialization Vector} (IV). The IV is a fresh value used together
with the key to generate a different keystream for each message.

The IV does not need to be secret, but the same IV must never be reused
with the same key.

\paragraph{Block Ciphers}

A block cipher encrypts fixed-size blocks of plaintext. If the plaintext
is shorter than one block, padding must be added.

Common examples are:

\begin{itemize}
    \item DES\@: 64-bit blocks, 56-bit key;
    \item 3DES\@: 64-bit blocks, 168-bit key (see LELEC2760 why 2DES doesn't exist ;-);
    \item AES\@: 128-bit blocks, keys of 128, 192, or 256 bits.
\end{itemize}

Block ciphers are not directly used on arbitrary-length messages. They
must be used with a \emph{mode of operation} which defines how a block cipher is applied to messages larger than one block. Examples include ECB\@, CBC\@, CFB\@, OFB\@, and CTR\@.


\paragraph{Cipher Block Chaining}

In Cipher Block Chaining (CBC), each plaintext block is XORed with the
previous ciphertext block then encrypted. This makes identical
plaintext blocks encrypt to different ciphertext blocks.

For the first block, there is no previous ciphertext block, so an
Initialization Vector is used. As with stream ciphers, the IV must be
fresh for each message.


\paragraph{Iterated Block Ciphers}

Many block ciphers are built by repeatedly applying a round function.
The original key is expanded into several round keys, and each round
transforms the internal state.

For example:

\begin{itemize}
    \item DES uses 16 rounds;
    \item 3DES applies DES three times;
    \item AES-128 uses 10 rounds.
\end{itemize}

\paragraph{DES}

DES is an old block cipher with a 56-bit key and 64-bit blocks. It uses
a Feistel structure, where each block is split into two halves $L_i$ and
$R_i$:

\[
R_i = f_i(R_{i-1}) \oplus L_{i-1}
\]

\[
L_i = R_{i-1}
\]

Feistel networks are useful because they are invertible, even if the
round function itself is not. Decryption is done by going through the network with the round keys in reverse order.

DES is no longer secure because its 56-bit key is too short and can be
brute-forced.

\paragraph{3DES}

3DES improves DES by applying DES three times with different keys:

\[
E_{3DES}((k_1,k_2,k_3),m)
=
E(k_1, D(k_2, E(k_3,m)))
\]

It is much stronger than DES, but also slower. Although it uses three
56-bit keys, its effective security is about 112 bits because of
meet-in-the-middle attacks.

\paragraph{AES}

AES is the modern standard block cipher. It has a block size of 128 bits and supports three key sizes:

\begin{itemize}
    \item AES-128: 10 rounds;
    \item AES-192: 12 rounds;
    \item AES-256: 14 rounds.
\end{itemize}


% \paragraph{Stream Ciphers vs Block Ciphers}

% Stream ciphers encrypt data symbol by symbol and can start immediately,
% without waiting for a full block. However, modifying individual parts of
% the ciphertext can sometimes have predictable effects.

% Block ciphers encrypt data block by block. They are often more robust
% against local modifications, but they require padding for short messages
% and errors can affect an entire block or more, depending on the mode.

\paragraph{Key Management}

A major problem in symmetric cryptography is key distribution. If every
pair of users needs a different shared key, then a system with $n$ users
requires:

\[
O(n^2)
\]

keys.

A possible solution is to use a \emph{Trusted Third Party} (TTP). Each
user shares one key with the TTP, and the TTP helps users establish
session keys for communication.

\paragraph{Replay Attacks}

A protocol can still be insecure even if encryption is used. For example,
an attacker may record an encrypted message and replay it later.

To prevent replay attacks, protocols usually include nonces, timestamps,
or sequence numbers, so that old messages cannot be reused successfully.

\subsection{Public-Key Cryptography}

In public-key cryptography, each user has a pair of keys: a public key
$pk$ and a private key $sk$. The public key can be shared with everyone,
while the private key must remain secret. To send a confidential message
to Bob, Alice encrypts the message using Bob's public key, and only Bob
can decrypt it using his private key.

\[
E(pk_{Bob}, m) = c
\]

\[
D(sk_{Bob}, c) = m
\]

A public-key cryptosystem consists of three algorithms:

\begin{itemize}
    \item $G$: generates a key pair $(pk, sk)$,
    \item $E(pk,m)$: encrypts a message using the public key,
    \item $D(sk,c)$: decrypts a ciphertext using the private key.
\end{itemize}

Correctness requires that decrypting an encrypted message gives back the
original message:

\[
D(sk, E(pk,m)) = m
\]

\paragraph{Advantages and Drawbacks}

The main advantage is key management: each user only needs one key pair
to communicate with many other users. However, public-key cryptography is
much slower than symmetric-key cryptography and usually relies on more
complex mathematical assumptions.

\paragraph{RSA}

RSA is a classical public-key cryptosystem. Its security relies on the
difficulty of deriving the private key from the public key, which is
related to the difficulty of factoring a large number $n$.

RSA key generation works as follows:

\begin{enumerate}
    \item choose two large primes $p$ and $q$,
    \item compute $n = p \cdot q$,
    \item compute $\varphi(n) = (p-1)(q-1)$,
    \item choose $e$ such that $\gcd(e,\varphi(n)) = 1$,
    \item compute $d$ such that $d \cdot e \equiv 1 \pmod{\varphi(n)}$.
\end{enumerate}

The public key is $(e,n)$ and the private key is $(d,n)$. The values
$p$, $q$, and $\varphi(n)$ must remain secret or be deleted after key
generation. Otherwise one could compute $d$ from $\varphi(n)$ and $e$.

Encryption and decryption are defined as:

\[
E((e,n),m) = m^e \bmod n
\]

\[
D((d,n),c) = c^d \bmod n
\]

The reason this works is that $d$ is chosen as the inverse of $e$ modulo
$\varphi(n)$, so decrypting reverses the encryption operation.

\paragraph{Security of Textbook RSA}

RSA in this simple form should not be used directly. It is deterministic:
the same message always produces the same ciphertext. This allows an
attacker to recognize repeated messages.

It also has algebraic structure. In particular, textbook RSA is
multiplicatively homomorphic:

\[
E(pk, m_1 \cdot m_2) = E(pk,m_1) \cdot E(pk,m_2)
\]

This can leak information and enable attacks. Practical RSA therefore
uses padding and randomness.

\paragraph{Hybrid Encryption}

In practice, RSA is not used to encrypt long messages directly because
it is too slow and can only encrypt messages smaller than the modulus
$n$. Instead, RSA is usually combined with symmetric cryptography.

A typical hybrid scheme works as follows:

\begin{enumerate}
    \item generate a random secret value $x$,
    \item encrypt it with RSA\@: $y = E_{RSA}(pk,x)$,
    \item derive a symmetric key $k = H(x)$,
    \item encrypt the message with a symmetric cipher:
    \[
    c = E_{sym}(k,m)
    \]
    \item send $(y,c)$.
\end{enumerate}

The receiver decrypts $y$ with their private key to recover $x$, derives
the same symmetric key $k = H(x)$, and then decrypts $c$.

\subsection{Diffie-Hellman Key Exchange}

Diffie-Hellman (DH) is a protocol that allows two parties to agree on a
shared secret over an insecure channel. This shared secret can then be
used as a key for symmetric encryption.

The protocol works as follows:

\begin{enumerate}
    \item Alice and Bob agree on a large prime $p$ and a generator $g$ (not needed to be private).
    \item Alice chooses a secret value $a$ and sends:
    \[
    X = g^a \bmod p
    \]
    \item Bob chooses a secret value $b$ and sends:
    \[
    Y = g^b \bmod p
    \]
    \item Alice computes:
    \[
    Y^a \bmod p = g^{ab} \bmod p
    \]
    \item Bob computes:
    \[
    X^b \bmod p = g^{ab} \bmod p
    \]
\end{enumerate}

Both parties obtain the same shared secret $g^{ab} \bmod p$, while an
eavesdropper only sees $p$, $g$, $g^a \bmod p$ and $g^b \bmod p$.

\paragraph{Security of Diffie-Hellman}

The security of DH is based on the difficulty of the discrete logarithm
problem. Computing $g^a \bmod p$ is easy, but recovering $a$ from
$g$, $p$ and $g^a \bmod p$ is hard when $p$ is large.

However, it is vulnerable to Man-in-the-Middle attacks where an attacker can intercept $g^a$ and $g^b$ and send both Alice and Bob $g^c$ each time decrypting and re-encrypting.

If a fresh DH exchange is performed for every connection, the protocol
provides \emph{forward secrecy}: even if a long-term secret is leaked
later, previously recorded sessions remain protected.

\paragraph{Elliptic Curve Diffie-Hellman}

Elliptic Curve Diffie-Hellman (ECDH) is a variant of DH based on the
elliptic curve discrete logarithm problem. Instead of working with
integers modulo $p$, it works with points on an elliptic curve.

ECDH provides similar security with much shorter keys. 

\subsection{Cryptographic Hash Functions}

A hash function maps an input of arbitrary length to a fixed-size output:

\[
H : {\{0,1\}}^* \rightarrow {\{0,1\}}^n
\]

A cryptographic hash function should be hard to invert and should behave
like a random function. The output is often called a \emph{message
digest}.

Examples include SHA-256 and SHA-512. Older functions such as MD5 and
SHA-1 are no longer recommended because practical attacks have been
found.

\paragraph{Desired Properties}

A cryptographic hash function should satisfy the following properties:

\begin{itemize}
    \item \textbf{Preimage resistance:} given $y$, it should be hard to
    find $m$ such that $H(m)=y$.
    \item \textbf{Collision resistance:} it should be hard to find two
    different messages $m_1$ and $m_2$ such that:
    \[
    H(m_1)=H(m_2)
    \]
    \item \textbf{Random-looking output:} the hash value should be
    indistinguishable from a random $n$-bit string.
\end{itemize}

\paragraph{Birthday Paradox}

For an $n$-bit hash function, there are $2^n$ possible outputs. Because
of the birthday paradox, a collision can be found after about
$2^{n/2}$ trials, not $2^n$.

For example, SHA-1 produces 160-bit hashes, so generic collision search
would take about $2^{80}$ trials. Since stronger attacks exist, SHA-1 is
not recommended anymore.

\paragraph{SHA-256}

% SHA-256 produces a 256-bit hash. Like many hash functions, it is based on
% the Merkle-Damg{\aa}rd construction: the message is split into blocks,
% and a compression function is applied repeatedly to update an internal
% state.

\subsection{Message Authentication Codes}

A Message Authentication Code (MAC) is used to verify the integrity and
authenticity of a message using a shared secret key.

A MAC consists of:

\[
S : K \times M \rightarrow T
\]

\[
V : K \times M \times T \rightarrow \{\text{yes},\text{no}\}
\]

The sender computes a tag:

\[
t = S(k,m)
\]

and sends $(m,t)$. The receiver checks the tag using the same key. If the
verification fails, the message is rejected.

\paragraph{Building MACs from Hash Functions}

A naive idea would be to build a MAC as:

\[
S(k,m)=H(m || k)
\]

or:

\[
S(k,m)=H(k || m)
\]

However, these constructions are insecure with common hash functions
based on Merkle-Damg{\aa}rd. In particular, $H(k||m)$ may be vulnerable
to length-extension attacks, where an attacker can extend the message
and compute a valid tag without knowing the key.

\paragraph{HMAC}

HMAC is the standard way to build a MAC from a hash function. Its basic
idea is to hide the intermediate hash value from the attacker:

\[
S(k,m)=H(k || H(k || m))
\]

The real construction is slightly more complex, but the intuition is that
the key is mixed into the hash computation in a safer way. A common
example is HMAC-SHA256.

